% GENERATED FILE. Edit canonical lessons, then run npm run generate:books.
% canonical-source-hash: fnv1a64:126dceb5e4a193e5
% canonical-lessons: ZH-C440-ganji, ZH-C440-shancun, ZH-W440-cun-village, ZH-C440-cunzhang, ZH-C440-cunmin, ZH-C440-jiesu

\chapter{A Mountain Village}
\label{ch:zh-a-mountain-village}
\hlchaptermodality{440}

\begin{chapteropening}
\textbf{By the end of this chapter:} \emph{I can talk about market day, a mountain village, its head and its villagers, and staying with a family.}

Complete the last lesson of chapter 440: I can talk about market day, a mountain village, its head and its villagers, and staying with a family.
\end{chapteropening}

\section[gǎnjí]{gǎnjí — to go to market day; to go to a country market}

\label{lesson:ZH-C440-ganji}



\begin{quote}
Before the new one: say the Mandarin for a duck, then the Mandarin for to take one's mind off things.
\end{quote}

\subsection*{You'll want to know: gǎnjí}
\emph{gǎnjí} — \textquotedblleft{}to go to market day; to go to a country market\textquotedblright{}.

\begin{quote}
\emph{nǎinai měi ge xīngqī qù gǎnjí} — Grandma goes to market day every week.
\end{quote}

Its characters come later; for now it is a word for the ear.

\begin{cousinweb}
\emph{gǎnjí} is built from \emph{gǎn} 'to hurry to' and \emph{jí} 'a gathering'. The second is inherited from Middle Chinese \emph{dzip} (Guangyun, in Baxter's transcription), glossed 'to gather'; no dictionary we use records an older form of \emph{gǎn} in this sense, so its history is not settled.
\end{cousinweb}

\begin{grammarlens}[title={What you've built}]
One more word for getting around.
\end{grammarlens}

\subsection*{Your turn}
\begin{itemize}
\raggedright
  \item \emph{Say it:} \emph{gǎnjí}
  \item \emph{Say it:} \emph{gǎnjí}, once more
  \item \emph{Say it:} \emph{nǎinai měi ge xīngqī qù gǎnjí} — Grandma goes to market day every week.
  \item \emph{Recall:} say the Mandarin for a duck, then the Mandarin for to take one's mind off things, then say \emph{gǎnjí} again
\end{itemize}

\begin{tcolorbox}[breakable,colback=teal!4,colframe=teal!35!black,title={Before you move on}]
What does gǎnjí mean? (\textquotedblleft{}To go to market day; to go to a country market\textquotedblright{}.) Say it once more.
\end{tcolorbox}
\section[\texorpdfstring{\zh{山村} (shāncūn)}{山村 (shāncūn)}]{\zh{山村} (shāncūn) — mountain village}

\label{lesson:ZH-C440-shancun}



\begin{quote}
Before the new one: say the Mandarin for to take one's mind off things, then the Mandarin for to go to market day.
\end{quote}

\subsection*{You'll want to know: \zh{山村}}
\textbf{\zh{山村}} — \emph{shāncūn} — \textquotedblleft{}mountain village\textquotedblright{}.

\begin{quote}
\emph{tā zhù zài yí ge xiǎo shāncūn} — He lives in a little mountain village.
\end{quote}

The new character, in pieces: \zh{村} — seven strokes in two pieces: on the left the tree piece of \zh{机}, four strokes; on the right the right piece of \zh{对}, three strokes. The next lesson puts it in your hand.

\begin{cousinweb}
\textbf{\zh{山村}} (\emph{shāncūn}) is built from \zh{山} 'a mountain' and \zh{村} 'a village', both inherited from Middle Chinese, \emph{srean} and \emph{tshwon} (Guangyun, in Baxter's transcription), which glosses the second as a farmstead.
\end{cousinweb}

\begin{grammarlens}[title={What you've built}]
One more word for getting around.
\end{grammarlens}

\subsection*{Your turn}
\begin{itemize}
\raggedright
  \item \emph{Say it:} \emph{shāncūn}
  \item \emph{Say it:} \emph{shāncūn}, once more
  \item \emph{Say it:} \emph{tā zhù zài yí ge xiǎo shāncūn} — He lives in a little mountain village.
  \item \emph{Recall:} say the Mandarin for to take one's mind off things, then the Mandarin for to go to market day, then say \emph{shāncūn} again
\end{itemize}

\begin{tcolorbox}[breakable,colback=teal!4,colframe=teal!35!black,title={Before you move on}]
What does \zh{山村} mean? (\textquotedblleft{}Mountain village\textquotedblright{}.) Say it once more.
\end{tcolorbox}
\section[\texorpdfstring{\zh{村} (cūn)}{村 (cūn)}]{\zh{村} — the tree piece, then the right piece of \zh{对}}

\label{lesson:ZH-W440-cun-village}



\begin{quote}
The word first, then the character inside it.

\begin{itemize}
\raggedright
  \item \emph{Recall:} say \emph{shāncūn} — \textbf{R1}, one lesson back, before it decays
\end{itemize}

\begin{quote}
\zh{机} \zh{对}
\end{quote}

\emph{Read it:} each one aloud, with nothing said first — \emph{jī}, \emph{duì}
\end{quote}

\subsection*{Script — the tree piece, then the right piece of \zh{对}}
\hlblockfigure{\detokenize{figures/ZH-W440-cun-village-filmstrip.pdf}}{How \zh{村} is written, stroke by stroke}

\begin{quote}
\zh{村}
\end{quote}

\textbf{Two pieces, left before right.} Count the pieces, not the strokes.

\textbf{The left piece, the tree piece, four strokes}, as on the left of \zh{机}: a bar, a vertical, a stroke falling to the left, and a short dot to the right.

\textbf{The right piece, three strokes}, as on the right of \zh{对}: a bar, a vertical with a hook at the foot, and a dot to the left of it.

4 + 3 = seven.

\subsection*{Your turn}
\begin{itemize}
\raggedright
  \item \emph{Write it:} \zh{村}, all seven, with the model in front of you
  \item \emph{Cover:} write it again from memory, and count the pieces, not the strokes
\end{itemize}

\begin{tcolorbox}[breakable,colback=teal!4,colframe=teal!35!black,title={Before you move on}]
How many strokes in \zh{村}? (\textbf{Seven}.) What is the left piece? (\textbf{The tree piece, as in \zh{机}}.) Where does the last dot go? (\textbf{To the left of the hooked vertical}.)

Source: \href{https://raw.githubusercontent.com/chanind/hanzi-writer-data/68d10a4b21150cae5e1ebbd223eed289cf32d90c/data/\%E6\%9D\%91.json}{Hanzi Writer Data \zh{村}, pinned stroke paths and medians}.
\end{tcolorbox}
\section[\texorpdfstring{\zh{村长} (cūnzhǎng)}{村长 (cūnzhǎng)}]{\zh{村长} (cūnzhǎng) — village head}

\label{lesson:ZH-C440-cunzhang}



\begin{quote}
Before the new one: say the Mandarin for to go to market day, then the Mandarin for a mountain village.
\end{quote}

\subsection*{You'll want to know: \zh{村长}}
\textbf{\zh{村长}} — \emph{cūnzhǎng} — \textquotedblleft{}village head\textquotedblright{}.

\begin{quote}
\emph{cūnzhǎng dài wǒmen qù tā jiā} — The village head takes us to his house.
\end{quote}

\begin{cousinweb}
\textbf{\zh{村长}} (\emph{cūnzhǎng}) is built from \zh{村} 'a village' and \zh{长} 'the head', both inherited from Middle Chinese, \emph{tshwon} and \emph{trjangX} (Guangyun, in Baxter's transcription), which glosses the second 'big'; in this sense \zh{长} is said zhǎng, not cháng.
\end{cousinweb}

\begin{grammarlens}[title={What you've built}]
One more word for getting around.
\end{grammarlens}

\subsection*{Your turn}
\begin{itemize}
\raggedright
  \item \emph{Say it:} \emph{cūnzhǎng}
  \item \emph{Say it:} \emph{cūnzhǎng}, once more
  \item \emph{Say it:} \emph{cūnzhǎng dài wǒmen qù tā jiā} — The village head takes us to his house.
  \item \emph{Recall:} say the Mandarin for to go to market day, then the Mandarin for a mountain village, then say \emph{cūnzhǎng} again
\end{itemize}

\begin{tcolorbox}[breakable,colback=teal!4,colframe=teal!35!black,title={Before you move on}]
What does \zh{村长} mean? (\textquotedblleft{}Village head\textquotedblright{}.) Say it once more.
\end{tcolorbox}
\section[cūnmín]{cūnmín — villager}

\label{lesson:ZH-C440-cunmin}



\begin{quote}
Before the new one: say the Mandarin for a mountain village, then the Mandarin for the village head.
\end{quote}

\subsection*{You'll want to know: cūnmín}
\emph{cūnmín} — \textquotedblleft{}villager\textquotedblright{}.

\begin{quote}
\emph{zhèr de cūnmín dōu hěn rèqíng} — The villagers here are all very friendly.
\end{quote}

Its characters come later; for now it is a word for the ear.

\begin{cousinweb}
\emph{cūnmín} is built from \zh{村} 'a village' and \emph{mín} 'the people', both inherited from Middle Chinese, \emph{tshwon} and \emph{mjin} (Guangyun, in Baxter's transcription).
\end{cousinweb}

\begin{grammarlens}[title={What you've built}]
One more word for getting around.
\end{grammarlens}

\subsection*{Your turn}
\begin{itemize}
\raggedright
  \item \emph{Say it:} \emph{cūnmín}
  \item \emph{Say it:} \emph{cūnmín}, once more
  \item \emph{Say it:} \emph{zhèr de cūnmín dōu hěn rèqíng} — The villagers here are all very friendly.
  \item \emph{Recall:} say the Mandarin for a mountain village, then the Mandarin for the village head, then say \emph{cūnmín} again
\end{itemize}

\begin{tcolorbox}[breakable,colback=teal!4,colframe=teal!35!black,title={Before you move on}]
What does cūnmín mean? (\textquotedblleft{}Villager\textquotedblright{}.) Say it once more.
\end{tcolorbox}
\section[jièsù]{jièsù — to lodge with a local family; to stay the night at someone's home}

\label{lesson:ZH-C440-jiesu}



\begin{quote}
Before the new one: say the Mandarin for the village head, then the Mandarin for a villager.
\end{quote}

\subsection*{You'll want to know: jièsù}
\emph{jièsù} — \textquotedblleft{}to lodge with a local family; to stay the night at someone's home\textquotedblright{}.

\begin{quote}
\emph{zhǎo bu dào jiǔdiàn de shíhou wǒmen zài cūnmín jiā jièsù} — When we can't find a hotel, we stay the night with a village family.
\end{quote}

Its characters come later; for now it is a word for the ear.

\begin{cousinweb}
\emph{jièsù} is built from \emph{jiè} 'to borrow' and \emph{sù} 'to stay the night', both inherited from Middle Chinese, \emph{tsjaeH} and \emph{sjuwk} (Guangyun, in Baxter's transcription), which glosses the first 'to borrow' and the second 'a lodging'.
\end{cousinweb}

\begin{grammarlens}[title={What you've built}]
One more word for getting around.
\end{grammarlens}

\subsection*{Your turn}
\begin{itemize}
\raggedright
  \item \emph{Say it:} \emph{jièsù}
  \item \emph{Say it:} \emph{jièsù}, once more
  \item \emph{Say it:} \emph{zhǎo bu dào jiǔdiàn de shíhou wǒmen zài cūnmín jiā jièsù} — When we can't find a hotel, we stay the night with a village family.
  \item \emph{Recall:} say the Mandarin for the village head, then the Mandarin for a villager, then say \emph{jièsù} again
\end{itemize}

\begin{tcolorbox}[breakable,colback=teal!4,colframe=teal!35!black,title={Before you move on}]
What does jièsù mean? (\textquotedblleft{}To lodge with a local family; to stay the night at someone's home\textquotedblright{}.) Say it once more.
\end{tcolorbox}
